\documentclass[11pt]{article}

% ── Packages ──────────────────────────────────────────────────────────────────
\usepackage{amsmath, amsthm, amssymb}
\usepackage{geometry}
\usepackage{graphicx}
\usepackage{tikz}
\usepackage{hyperref}
\usepackage{enumitem}
\usepackage{xcolor}
\usepackage{mdframed}
\usepackage{float}
\usepackage{algorithm}
\usepackage{algpseudocode}
\usepackage{algpseudocode}

\geometry{margin=1in}
\hypersetup{colorlinks=true, linkcolor=blue!60!black, urlcolor=blue!60!black}

% ── Theorem Environments ──────────────────────────────────────────────────────
\theoremstyle{plain}
\newtheorem{theorem}{Theorem}[section]
\newtheorem{lemma}[theorem]{Lemma}
\newtheorem{proposition}[theorem]{Proposition}

\theoremstyle{definition}
\newtheorem{definition}[theorem]{Definition}
\newtheorem{remark}[theorem]{Remark}

% ── Macros ────────────────────────────────────────────────────────────────────
\newcommand{\Z}{\mathbb{Z}}
\newcommand{\R}{\mathbb{R}}
\newcommand{\norm}[1]{\left\lVert #1 \right\rVert}
\newcommand{\floor}[1]{\left\lfloor #1 \right\rfloor}


% ── Highlight box ─────────────────────────────────────────────────────────────
\newmdenv[
  backgroundcolor=blue!5,
  linecolor=blue!40,
  linewidth=1.2pt,
  roundcorner=4pt,
  skipabove=8pt,
  skipbelow=8pt
]{keybox}

\begin{document}

\noindent
\begin{tabular}{|p{0.72\textwidth} r|}
\hline
\textbf{Recreational Mathematics Club} & \textit{Discussion \#4} \\[4pt]
\multicolumn{2}{|c|}{\Large Linear Congruential Generators and Lattice Attacks} \\[4pt]
Speaker: \textit{Dhairya Baxi, Mahek Shamsukha} & Indian Institute of Science \\
\hline
\end{tabular}

% ══════════════════════════════════════════════════════════════════════════════
\section{Linear Congruential Generators}
% ══════════════════════════════════════════════════════════════════════════════

\begin{definition}
A \emph{linear congruential generator} (LCG) with parameters $a, b, m \in \Z$
produces a sequence via the recurrence
\[
  x_{i+1} \equiv a x_i + b \pmod{m}.
\]
Starting from a \emph{seed} $x_0$, the full sequence $x_0, x_1, x_2, \ldots$
is determined. Each $x_i \in \{0, 1, \ldots, m-1\}$.
\end{definition}

LCGs are among the oldest and most widely used pseudorandom number generators,
appearing in standard libraries of many programming languages.

% ══════════════════════════════════════════════════════════════════════════════
\section{Refresher: GCD and Extended GCD}
% ══════════════════════════════════════════════════════════════════════════════

\subsection{Greatest Common Divisor}

The \emph{greatest common divisor} $\gcd(a,b)$ is the largest integer that
divides both $a$ and $b$.

The key identity is:
\[
  \gcd(a,b) = \gcd(b,\, a \bmod b).
\]

This leads to a simple iterative algorithm:

\begin{algorithm}[H]
\caption{\texttt{gcd(a,b)}}
\begin{algorithmic}[1]
    \State \textbf{Input:} Integers $a,b \geq 0$
    \State \textbf{Output:} $\gcd(a,b)$
    \While{$b \neq 0$}
        \State $(a,b) \gets (b,\ a \bmod b)$
    \EndWhile
    \State \Return $a$
\end{algorithmic}
\end{algorithm}

\begin{keybox}
\textbf{Idea.}
Each step replaces $(a,b)$ with a smaller pair having the same gcd.
Thus the algorithm quickly converges.
\end{keybox}

\subsection{Extended GCD}

Beyond computing the gcd, we often want integers $x,y$ such that
\[
  ax + by = \gcd(a,b).
\]

\begin{theorem}[Bézout]
For any integers $a,b$, there exist $x,y \in \Z$ such that
\[
  ax + by = \gcd(a,b).
\]
\end{theorem}

The extended Euclidean algorithm computes these coefficients.

\begin{algorithm}[H]
\caption{\texttt{extended\_gcd(a,b)}}
\begin{algorithmic}[1]
    \State \textbf{Input:} Integers $a,b \geq 0$
    \State \textbf{Output:} $(g,x,y)$ such that $ax + by = g = \gcd(a,b)$
    \If{$b = 0$}
        \State \Return $(a,\ 1,\ 0)$
    \EndIf
    \State $(g,x,y) \gets \texttt{extended\_gcd}(b,\ a \bmod b)$
    \State \Return $(g,\ y,\ x - \floor{a/b} \cdot y)$
\end{algorithmic}
\end{algorithm}

\begin{keybox}
\textbf{Idea.}
We recursively reduce the problem using
\[
  \gcd(a,b) = \gcd(b, a \bmod b),
\]
and then reconstruct the coefficients $(x,y)$ on the way back.
\end{keybox}

\subsection{Complexity}

Both \texttt{gcd} and \texttt{extended\_gcd} run in
\[
  O(\log \min(a,b))
\]
time.

% ══════════════════════════════════════════════════════════════════════════════
\section{Multiplicative Inverse via Extended GCD}
% ══════════════════════════════════════════════════════════════════════════════

We say $a^{-1}$ is a \emph{multiplicative inverse of $a$ modulo $m$} if
\[
  a \cdot a^{-1} \equiv 1 \pmod{m}.
\]

\begin{proposition}
The inverse $a^{-1} \pmod{m}$ exists if and only if $\gcd(a, m) = 1$.
\end{proposition}

\begin{proof}
$(\Leftarrow)$
Suppose $\gcd(a,m) = 1$. By B\'{e}zout, there exist integers $s, t$ with
$sa + tm = 1$. Reducing modulo $m$ gives $sa \equiv 1 \pmod{m}$,
so $s \bmod m$ is the desired inverse.

$(\Rightarrow)$
We prove the contrapositive: if $\gcd(a,m) = d > 1$, then no inverse exists.
Suppose for contradiction that some $u$ satisfies $au \equiv 1 \pmod{m}$,
i.e.\ $au = 1 + qm$ for some integer $q$, or equivalently $au - qm = 1$.
But $d \mid a$ and $d \mid m$, so $d \mid (au - qm) = 1$,
which is impossible for $d > 1$.
\end{proof}

\begin{keybox}
\textbf{Algorithm.}
Run the extended Euclidean algorithm on $(a, m)$.
\begin{itemize}
  \item If $\gcd(a,m) = 1$: return $s \bmod m$ as the inverse.
  \item If $\gcd(a,m) \neq 1$: report that no inverse exists.
\end{itemize}
\end{keybox}

% ══════════════════════════════════════════════════════════════════════════════
\section{The \texttt{next} and \texttt{forward} Routines}
% ══════════════════════════════════════════════════════════════════════════════

\subsection{One Step Forward}

\[
  \texttt{next}(x) := ax + b \bmod m.
\]

\subsection{$n$ Steps Forward}

\textbf{Naive approach:} Apply \texttt{next} repeatedly, costing $O(n)$ time.

\textbf{Efficient closed form.} Unrolling the recurrence gives:
\[
  \texttt{forward}(x, n) := a^n x + b \cdot \frac{a^n - 1}{a - 1} \pmod{m}
  \quad (a \neq 1).
\]

For $a = 1$ the generator is simply $x_n = x_0 + nb \bmod m$.

\begin{keybox}
\textbf{Efficiency.}
Both $a^n \bmod m$ and $(a^n - 1)(a-1)^{-1} \bmod m$ can be computed
in $O(\log n)$ time via fast modular exponentiation and the extended GCD
algorithm. Thus \texttt{forward} runs in $O(\log n)$ time regardless of $n$.
\end{keybox}

% ══════════════════════════════════════════════════════════════════════════════
\section{The \texttt{prev} and \texttt{backward} Routines}
% ══════════════════════════════════════════════════════════════════════════════

\subsection{One Step Backward}

Provided $\gcd(a, m) = 1$ so that $a^{-1}$ exists:
\[
  \texttt{prev}(x) := a^{-1} x - a^{-1} b \pmod{m}
  = a^{-1}(x - b) \pmod{m}.
\]

\subsection{$n$ Steps Backward}

\[
  \texttt{backward}(x, n) := \texttt{prev}^n(x).
\]

\begin{keybox}
\textbf{Efficient computation.}
By an identical argument to \texttt{forward}, \texttt{backward}$(x, n)$
admits a closed form with $a^{-n} \bmod m$, computable in $O(\log n)$ time.
Explicitly:
\[
  \texttt{backward}(x, n) = (a^{-1})^n x + b \cdot \frac{(a^{-1})^n - 1}{a^{-1} - 1} \pmod m.
\]
\end{keybox}

% ══════════════════════════════════════════════════════════════════════════════
\section{Problem Formulation}
% ══════════════════════════════════════════════════════════════════════════════

Fix parameters:
\begin{itemize}
  \item $n \geq 0$: the index of the first constrained output,
  \item $k \geq 1$: number of constrained outputs,
  \item $m_i, M_i \in [0, m)$ for $i \in [k]$: desired bounds on each output.
\end{itemize}

\begin{keybox}
\textbf{Problem.}
Find a seed $x_1 \in \{0, \ldots, m-1\}$ such that, with $x_{i+1} = \texttt{next}(x_i)$,
\[
  x_{n+i} \in [m_i,\, M_i] \quad \text{for all } i \in [k].
\]
\end{keybox}

% ══════════════════════════════════════════════════════════════════════════════
\section{Reduction to an Integer Linear Program}
% ══════════════════════════════════════════════════════════════════════════════

\subsection{Reduction to $n = 0$}

The constraint window starts at $x_{n+1}$. Using \texttt{forward}, the constrained
outputs are $x_{n+1}, x_{n+2}, \ldots, x_{n+k}$. Relabelling: set $y_1 = x_{n+1}$
and note that $y_1 = \texttt{forward}(x_1, n)$. Since \texttt{forward} and
\texttt{backward} are invertible (when $\gcd(a,m)=1$), finding a valid $x_1$
is equivalent to finding a valid $y_1$. We henceforth seek $y_1$ such that
\[
  y_{i} \in [m_i,\, M_i] \quad \text{for } i = 1, \ldots, k,
\]
where $y_{i+1} = \texttt{next}(y_i)$

\subsection{Variables}

Define:
\begin{itemize}
  \item $y_1 \in \{0,\ldots,m-1\}$: the unknown seed (after shifting by $n$),
  \item $s_i \in \Z$ for $i = 2, \ldots, k$: the integer such that
        $y_i + m s_i \in [0, m)$, i.e.\ the \emph{modular correction} ensuring
        $y_i$ is the true LCG output.
\end{itemize}

Concretely, removing the modulus from $y_{i+1} = a y_i + b \bmod m$ gives the
exact equation
\[
  y_{i+1} = a y_i + b - m s_{i+1},
\]
where $s_{i+1}$ is the quotient $\lfloor (a y_i + b)/m \rfloor$. These $s_i$
are not free: they are determined by $y_1$, but since they are integers they
serve as auxiliary integer variables in the ILP.

\subsection{Matrix Form}

Write $z = (y_1,\, s_2,\, s_3,\, \ldots,\, s_k)^\top \in \Z^k$. Unrolling
the recurrence, each constrained output $y_i$ is an affine function of $z$:
\[
  \begin{pmatrix} y_2 \\ y_3 \\ \vdots \\ y_{k+1} \end{pmatrix}
  = A z + \mathbf{b},
\]
for an explicit integer matrix $A$ and constant vector $\mathbf{b}$ depending
on $a, b, m$. The constraints $y_{1+i} \in [m_i, M_i]$ then read:
\[
  \mathbf{m} \leq A z + \mathbf{b} \leq \mathbf{M},
\]
where $\mathbf{m} = (m_1, \ldots, m_k)^\top$ and $\mathbf{M} = (M_1, \ldots, M_k)^\top$.

\subsection{Normalising the Bounds}

Subtracting $\mathbf{b}$ from all sides:
\[
  \mathbf{m} - \mathbf{b} \leq Az \leq \mathbf{M} - \mathbf{b}.
\]
The entries of $\mathbf{m} - \mathbf{b}$ and $\mathbf{M} - \mathbf{b}$ may
lie outside $[0,m)$. We bring them inside $[0,m)$ by adding appropriate
multiples of $m$ to each component. Adding $c_i \cdot m$ to the $i$-th
component of both sides is equivalent to adding $c_i \cdot m$ times the
$i$-th row of $A^{-1}$ to $z$, which can be absorbed into a change of
basis for $A$ — i.e.\ it corresponds to modifying the entries of $A$.
After this adjustment, define
\[
  \bar{\mathbf{m}} = \mathbf{m} - \mathbf{b} + C\mathbf{m}_{\text{shift}},
  \qquad
  \bar{\mathbf{M}} = \mathbf{M} - \mathbf{b} + C\mathbf{m}_{\text{shift}},
\]
and a modified matrix $\bar{A}$, so that the system becomes
\[
  \bar{\mathbf{m}} \leq \bar{A} z \leq \bar{\mathbf{M}},
  \quad \bar{\mathbf{m}},\, \bar{\mathbf{M}} \in [0,m)^k.
\]

\begin{keybox}
\textbf{Summary.}
The LCG constraint problem reduces to finding an integer vector
$z = (y_1, s_2, \ldots, s_k)^\top \in \Z^k$ satisfying
$\bar{\mathbf{m}} \leq \bar{A} z \leq \bar{\mathbf{M}}$,
with all bounds in $[0,m)^k$.
This is an \emph{integer linear program} (ILP).
\end{keybox}

% ══════════════════════════════════════════════════════════════════════════════
\section{Geometric and Lattice Interpretation}
% ══════════════════════════════════════════════════════════════════════════════

\subsection{Lattice Points in a Box}

The image of $\Z^k$ under $\bar{A}$ is the \emph{lattice}
$\mathcal{L}(\bar{A}) = \{ \bar{A} z : z \in \Z^k \}$.
Our problem asks: does $\mathcal{L}(\bar{A})$ contain any point inside
the axis-aligned box $[\bar{\mathbf{m}},\, \bar{\mathbf{M}}] \subset [0,m)^k$?

\subsection{Transforming to Integer Space}

Assuming $\bar{A}$ is square and invertible, applying $\bar{A}^{-1}$ to
the inequality $\bar{\mathbf{m}} \leq \bar{A} z \leq \bar{\mathbf{M}}$
transforms the problem: we need an integer point $z \in \Z^k$ lying inside
the \emph{parallelogram}
\[
  P = \bar{A}^{-1}[\bar{\mathbf{m}},\, \bar{\mathbf{M}}]
  = \{ z \in \R^k : \bar{\mathbf{m}} \leq \bar{A} z \leq \bar{\mathbf{M}} \}.
\]

\begin{keybox}
\textbf{Geometric reformulation.}
The axis-aligned box $[\bar{\mathbf{m}}, \bar{\mathbf{M}}]$ in output space
becomes a skewed parallelogram $P$ in the integer variable space of $z$.
The columns of $\bar{A}^{-1}$ determine how the coordinate axes are sheared.
We seek any $z \in \Z^k \cap P$.
\end{keybox}

% ══════════════════════════════════════════════════════════════════════════════
\section{Branch and Bound}
% ══════════════════════════════════════════════════════════════════════════════

A classical approach to ILPs is \emph{branch and bound}:

\begin{enumerate}
  \item \textbf{Relax} integrality: solve the LP relaxation over $\R^{k+1}$.
  \item If the solution is integral, we are done.
  \item Otherwise, pick a fractional variable $y_j \notin \Z$ and
        \textbf{branch}: create two subproblems with $y_j \leq \lfloor y_j^* \rfloor$
        and $y_j \geq \lceil y_j^* \rceil$ respectively.
  \item \textbf{Bound}: prune subproblems whose LP relaxation is infeasible
        or worse than the best solution found.
\end{enumerate}

\begin{keybox}
\textbf{Geometric issue.}
The efficiency of branch and bound depends critically on the shape of $P$.
If the parallelogram is highly \emph{skewed} — i.e., elongated along
directions misaligned with the integer lattice — the search tree becomes
exponentially large and the method is slow.
\end{keybox}

% ══════════════════════════════════════════════════════════════════════════════
\section{Basis Choice and Lattice Reduction}
% ══════════════════════════════════════════════════════════════════════════════

\subsection{Different Bases, Same Lattice}

A lattice $\mathcal{L}$ has infinitely many bases. If $B = (b_1, \ldots, b_n)$
is one basis and $U$ is any unimodular matrix (integer matrix with $\det U = \pm 1$),
then $BU$ is another valid basis generating the same lattice.

\begin{keybox}
\textbf{Observation.}
Two bases for the same lattice can produce parallelograms of wildly different
shapes. A skewed basis yields a thin, elongated parallelogram; an orthogonal
basis yields a compact, box-like one.
\end{keybox}

\subsection{What Makes a Good Basis?}

We want a basis $B = (b_1, \ldots, b_n)$ where:
\begin{itemize}
  \item The vectors $b_i$ are \emph{reasonably short} (no unnecessarily long directions),
  \item The vectors are \emph{nearly orthogonal} (small off-diagonal
        Gram--Schmidt coefficients),
  \item The Gram--Schmidt lengths $\|b_i^*\|$ do not drop too rapidly
        (avoiding highly skewed configurations).
\end{itemize}

Such a basis makes the corresponding parallelepiped well-shaped and
better aligned with the coordinate axes, dramatically improving the
performance of branch-and-bound methods.

\begin{keybox}
\textbf{Motivation for the next step.}
Finding an \emph{optimal} basis (e.g.\ shortest or most orthogonal) is
NP-hard in general. However, efficient approximation algorithms exist.
The most celebrated is the \emph{LLL algorithm} (Lenstra--Lenstra--Lov\'{a}sz, 1982),
which computes, in polynomial time, a basis that is nearly orthogonal and
well-balanced. In particular, it guarantees that the first basis vector is
within an exponential factor of the shortest nonzero lattice vector.
\end{keybox}
% ══════════════════════════════════════════════════════════════════════════════
\section{The LLL Algorithm}
% ══════════════════════════════════════════════════════════════════════════════

\subsection{$\delta$-LLL Reduced Basis}

Let $B = (b_1, \dots, b_n)$ be a basis of a lattice $\mathcal{L} \subset \R^m$.
Let $b_1^*, \dots, b_n^*$ be its Gram--Schmidt orthogonalization:
\[
  b_i = b_i^* + \sum_{j<i} \mu_{i,j} b_j^*.
\]

\begin{definition}
Let $\delta \in (1/4,1)$. The basis $B$ is called \emph{$\delta$-LLL reduced} if:
\begin{enumerate}
  \item For all $i > j$, $|\mu_{i,j}| \leq \tfrac{1}{2}$,
  \item For all $i = 1, \dots, n-1$,
  \[
    \|b_{i+1}^*\|^2 \;\geq\; (\delta - \mu_{i+1,i}^2)\,\|b_i^*\|^2.
  \]
\end{enumerate}
\end{definition}

% ══════════════════════════════════════════════════════════════════════════════
\subsection{Bound on the Shortest Vector}
% ══════════════════════════════════════════════════════════════════════════════

\begin{theorem}
If $B$ is a $\delta$-LLL reduced basis of $\mathcal{L}$, then
\[
  \|b_1\| \;\leq\; 2^{\frac{n-1}{2}} \cdot \lambda_1(\mathcal{L}),
\]
where $\lambda_1(\mathcal{L})$ is the length of the shortest nonzero lattice vector.
\end{theorem}

% ══════════════════════════════════════════════════════════════════════════════
\subsection{LLL Algorithm}
% ══════════════════════════════════════════════════════════════════════════════

\begin{algorithm}[H]
\caption{\texttt{LLL(B, $\delta$)}}
\begin{algorithmic}[1]
    \State \textbf{Input:} $B \in \Z^{m \times n}$, $\delta \in (1/4,1)$
    \State \textbf{Output:} A $\delta$-LLL reduced basis generating the same lattice
    \State $B \gets \textsc{SizeReduce}(B)$
    \If{ $\exists i$ such that $\|b_{i+1}^*\|^2 < (\delta - \mu_{i+1,i}^2)\,\|b_i^*\|^2$}
        \State Swap $b_i, b_{i+1}$
        \State \Return $\texttt{LLL}(B, \delta)$
    \EndIf
    \State \Return $B$
\end{algorithmic}
\end{algorithm}

% ══════════════════════════════════════════════════════════════════════════════
\subsection{Size Reduction via Column Operations}
% ══════════════════════════════════════════════════════════════════════════════

Recall the Gram--Schmidt decomposition
\[
  B = B^* M,
\]
where $B^* = (b_1^*, \dots, b_n^*)$ has orthogonal columns and
$M = (\mu_{i,j})$ is upper triangular with $\mu_{i,i} = 1$.

Size reduction corresponds to modifying the columns of $B$ using
\emph{integer column operations} so that
\[
  |\mu_{i,j}| \leq \tfrac{1}{2} \quad \text{for all } i > j.
\]

\medskip

We proceed column by column. For $j = 1, \dots, n$, and for
$i = j-1, j-2, \dots, 1$, perform the operation
\[
  b_j \;\leftarrow\; b_j - \lfloor \mu_{j,i} \rceil \, b_i,
\]
where $\lfloor \cdot \rceil$ denotes rounding to the nearest integer.

\medskip

This is an elementary column operation on $B$, and corresponds to updating
the matrix $M$ by
\[
  \mu_{j,i} \;\leftarrow\; \mu_{j,i} - \lfloor \mu_{j,i} \rceil,
\]
which ensures
\[
  |\mu_{j,i}| \leq \tfrac{1}{2}.
\]

\begin{keybox}
\textbf{Key point.}
Size reduction is achieved purely by integer column operations on $B$,
preserving the lattice, while ensuring all Gram--Schmidt coefficients
satisfy $|\mu_{i,j}| \leq \tfrac{1}{2}$.
\end{keybox}

% ══════════════════════════════════════════════════════════════════════════════
\subsection{Existence of LLL-Reduced Basis}
% ══════════════════════════════════════════════════════════════════════════════

\begin{theorem}
Let $B \in \Z^{m \times n}$ be an integer basis and let $\delta \in (1/4,1)$.
Then there exists an integer basis $B'$ such that:
\begin{enumerate}
  \item $B'$ is $\delta$-LLL reduced,
  \item $B'$ generates the same lattice as $B$.
\end{enumerate}
\end{theorem}

\begin{proof}
We prove correctness and finiteness of the LLL algorithm.

\medskip
\noindent\textbf{(Correctness).}
If the algorithm terminates, then no Lovász condition is violated, and the
basis is size-reduced. Hence the output basis is $\delta$-LLL reduced.

\medskip
\noindent\textbf{(Lattice preservation).}
All operations used are integer column operations:
\begin{itemize}
  \item size reduction replaces $b_j$ by $b_j - k b_i$,
  \item swaps exchange two basis vectors.
\end{itemize}
Thus the output basis is obtained from the input by unimodular transformations,
and therefore generates the same lattice.

\medskip
  \noindent\textbf{(Finiteness)}
Define a potential
\[
  \Phi(B)
  \;:=\;
  \prod_{i=1}^n \det\!\big(\mathcal{L}(b_1,\dots,b_i)\big),
\]
where $\mathcal{L}(b_1,\dots,b_i)$ denotes the lattice generated by
$b_1,\dots,b_i$.

Equivalently, since
\[
  \det\!\big(\mathcal{L}(b_1,\dots,b_i)\big)
  = \text{vol of the fundamental parallelepiped}
  = \prod_{j=1}^i \|b_j^*\|,
\]
we may write
\[
  \Phi(B)
  =
  \prod_{i=1}^n \prod_{j=1}^i \|b_j^*\|.
\]

\medskip
\noindent\emph{Effect of size reduction.}
Size reduction does not change the Gram--Schmidt vectors $b_j^*$.
Hence $\Phi(B)$ remains unchanged.

\medskip
\noindent\emph{Effect of swapping $b_i, b_{i+1}$.}
All prefixes except the $i$-th remain unchanged, so only the factor
\[
  \det\!\big(\mathcal{L}(b_1,\dots,b_i)\big)
\]
changes.

Let $\Phi'$ denote the potential after swapping $b_i$ and $b_{i+1}$.

First observe how the Gram--Schmidt vectors change. Consider the prefix
\[
  (b_1, \dots, b_{i-1}, b_{i+1}).
\]
For this sequence, the Gram--Schmidt vectors are
\[
  b_1^*, \dots, b_{i-1}^*, \; b_{i+1}^{*\prime},
\]
where the first $i-1$ vectors remain unchanged, and
\[
  b_{i+1}^{*\prime}
  =
  b_{i+1} - \sum_{j=1}^{i-1} \mu_{i+1,j} b_j^*
  =
  b_{i+1}^* + \mu_{i+1,i} b_i^*.
\]

Thus,
\[
  \|b_{i+1}^{*\prime}\|^2
  =
  \|b_{i+1}^*\|^2 + \mu_{i+1,i}^2 \|b_i^*\|^2.
\]

Now, since only the $i$-th prefix changes in the potential, we have
\[
  \frac{\Phi'}{\Phi}
  =
  \frac{\det(\mathcal{L}(b_1,\dots,b_{i-1},b_{i+1}))}
       {\det(\mathcal{L}(b_1,\dots,b_i))}
  =
  \frac{\|b_{i+1}^{*\prime}\|}{\|b_i^*\|}.
\]

If the Lovász condition is violated, then
\[
  \|b_{i+1}^*\|^2 < (\delta - \mu_{i+1,i}^2)\|b_i^*\|^2,
\]
and hence
\[
  \|b_{i+1}^{*\prime}\|^2
  <
  \delta \|b_i^*\|^2.
\]

Taking square roots,
\[
  \frac{\Phi'}{\Phi}
  <
  \sqrt{\delta}
  <
  1.
\]

\medskip
\noindent\emph{Finiteness.}
Observe that for each $i$,
\[
  \det\!\big(\mathcal{L}(b_1,\dots,b_i)\big)
  =
  \sqrt{\det(B_i^T B_i)},
\]
where $B_i$ is the matrix with columns $b_1,\dots,b_i$. Since $B_i$ has
integer entries, $\det(B_i^T B_i)$ is an integer, and hence each factor in
$\Phi(B)$ is a square root of a positive integer.

Therefore, $\Phi(B)$ itself is a product of square roots of integers, and
can take values only in a discrete subset of $\R_{>0}$. In particular,
in the interval $(0, \Phi(B^{\text{in}})]$, it can assume only finitely
many values.

Since each swap strictly decreases $\Phi(B)$ by at least a constant factor
$\sqrt{\delta} < 1$, the sequence of potentials is strictly decreasing and
cannot be infinite.

\medskip
\noindent
Therefore, the algorithm terminates.

\medskip
Combining the above, a $\delta$-LLL reduced basis $B'$ exists.
\end{proof}
% ══════════════════════════════════════════════════════════════════════════════
\section*{References}
% ══════════════════════════════════════════════════════════════════════════════

\begin{itemize}
  \item \textit{Minecraft 12-Eyed Seed Explanation (YouTube video)}\\
        \url{https://www.youtube.com/watch?v=mc9w2iD3Gzs}

  \item D.\ Micciancio, \textit{CSE206A: Lattice Algorithms and Applications
        (Winter 2010)}\\
        Lecture notes, UC San Diego\\
        \url{https://cseweb.ucsd.edu/classes/wi10/cse206a/}

  \item O.\ Regev, \textit{Lattices in Computer Science (Fall 2004)}\\
        Lecture notes, Tel Aviv University\\
        \url{https://cims.nyu.edu/~regev/teaching/lattices_fall_2004/}
\end{itemize}
\newpage
% ══════════════════════════════════════════════════════════════════════════════
\appendix
\section{Appendix: Minecraft 12-Eyed Seed}
% ══════════════════════════════════════════════════════════════════════════════

In \emph{Minecraft}, an End Portal consists of 12 slots, each of which
requires an \emph{Ender Eye} to be filled. When a world is generated,
each slot may already be pre-filled.

Internally, this is determined using calls to a pseudorandom generator
(\texttt{Random::nextFloat()} in Java), which maps the internal seed to a
value in $[0,1)$. A slot is filled if this value exceeds $0.9$.

\medskip

The Java random number generator is a linear congruential generator:
\[
  \text{seed}_{i+1}
  \equiv
  25214903917 \cdot \text{seed}_i + 11
  \pmod{2^{48}}.
\]

\medskip

In the world-generation pipeline, the first $760$ calls are irrelevant for
the portal. Thus, we begin at $\text{seed}_{761}$ and require that the next
$12$ calls all satisfy
\[
  \texttt{nextFloat()} > 0.9.
\]

This translates to the constraint
\[
  \text{seed}_i \in [253327472328704,\; 2^{48}-1],
  \qquad i = 761, \dots, 772,
\]
since
\[
  0.9 \cdot 2^{48} = 253327472328704.
\]

\begin{keybox}
\textbf{Connection to our problem.}
This is an instance of the LCG constraint problem with:
\begin{itemize}
  \item modulus $m = 2^{48}$,
  \item multiplier $a = 25214903917$, increment $b = 11$,
  \item $n = 760$ unconstrained steps,
  \item $k = 12$ constrained steps,
  \item uniform bounds $m_i = 253327472328704$ and $M_i = 2^{48}-1$.
\end{itemize}
\end{keybox}

Thus, finding a 12-eyed seed reduces to finding integer points satisfying
linear constraints derived from an LCG --- precisely the type of problem
we studied using lattice methods and the LLL algorithm.
\end{document}
